\chapter{Il tessuto connettivo}

% ---------------------------------------------------------------------------
\section{Classificazione}

Il \textbf{tessuto connettivo} si distingue dagli altri tessuti per l'abbondante
\textbf{matrice extracellulare} che separa le cellule tra loro. Si suddivide in
tre grandi categorie:

\begin{itemize}
  \item \textbf{Tessuto connettivo propriamente detto}
    \begin{itemize}
      \item \textit{Lasso}: le fibre formano una rete aperta; comprende il
        tessuto areolare, il tessuto adiposo e il tessuto reticolare.
      \item \textit{Denso}: le fibre sono strettamente impacchettate; comprende
        il denso regolare, il denso irregolare e il tessuto elastico.
    \end{itemize}
  \item \textbf{Tessuto connettivo liquido}
    \begin{itemize}
      \item \textit{Sangue}: contenuto nel sistema circolatorio.
      \item \textit{Linfa}: contenuta nel sistema linfatico.
    \end{itemize}
  \item \textbf{Tessuto connettivo di sostegno}
    \begin{itemize}
      \item \textit{Cartilagine}: matrice solida ed elastica; si distingue in
        ialina, elastica e fibrosa.
      \item \textit{Osso}: matrice solida a struttura cristallina.
    \end{itemize}
\end{itemize}

% ---------------------------------------------------------------------------
\section{Il sangue}

\subsection{Generalità e composizione}

Il \textbf{sangue} è un tessuto connettivo allo stato liquido rinchiuso in un
sistema di canali comunicanti (vasi arteriosi e venosi). Fornisce nutrienti
(vitamine, elettroliti) e O\textsubscript{2} alle cellule corporee e trasporta
i prodotti di scarto verso gli organi emuntori.

Centrifugando un campione di sangue si ottengono tre frazioni:

\begin{itemize}
  \item \textbf{Plasma} (55\,\%): parte liquida contenente acqua, proteine,
    lipidi, glucosio, aminoacidi e ioni.
  \item \textbf{Globuli bianchi e piastrine} ($\sim$1\,\%): strato intermedio
    detto \textit{buffy coat}.
  \item \textbf{Globuli rossi} (45\,\%): frazione più densa, detta
    \textit{ematocrito}.
\end{itemize}

Il volume di sangue in un individuo adulto è circa il 7--8\,\% del peso
corporeo (4,5--6 litri). Le \textbf{vene} fungono da tampone: consentono di
tollerare una perdita acuta fino al 20\,\% del volume totale
($\sim$1\,litro) senza gravi conseguenze. Il plasma è in costante equilibrio
osmotico con i liquidi extracellulari dei tessuti.

\subsection{Funzioni del sangue}

\begin{itemize}
  \item \textbf{Trasporto}: di sostanze nutritive, ormoni e gas disciolti.
  \item \textbf{Riserva}: di nutrienti e fattori della coagulazione.
  \item \textbf{Termoregolazione}: redistribuzione del calore corporeo.
  \item \textbf{Difesa}: contro patogeni e sostanze estranee.
  \item \textbf{Antiemorragica}: coagulazione e riparazione delle lesioni vascolari.
  \item \textbf{Omeostasi}: regolazione del pH e della composizione
    elettrolitica dei liquidi interstiziali.
\end{itemize}

% ---------------------------------------------------------------------------
\section{Il plasma e le proteine plasmatiche}

\subsection{Plasma e siero}

Il \textbf{plasma} differisce dal \textbf{siero} per la presenza dei fattori
della coagulazione: quando il fibrinogeno precipita formando fibrina, il plasma
si trasforma in siero.

\begin{itemize}
  \item \textbf{Siero} = albumina + globuline.
  \item \textbf{Plasma} = albumina + globuline + fibrinogeno + fattori della
    coagulazione.
\end{itemize}

\subsection{Proteine plasmatiche}

Le proteine plasmatiche sono catene polipeptidiche, talvolta coniugate con
lipidi o carboidrati (lipoproteine, glicoproteine). Svolgono funzioni di:
difesa immunitaria umorale, mantenimento della pressione oncotica, trasporto
di sostanze insolubili in acqua, regolazione del pH e coagulazione del sangue.

Sono sintetizzate principalmente dagli \textbf{epatociti} (albumina e
globuline), dalle \textbf{plasmacellule} (immunoglobuline) e
dall'\textbf{intestino} (lipoproteine). Il loro metabolismo è rapido: circa
il 9--15\,\% viene rinnovato ogni giorno.

\subsection{Frazioni elettroforetiche e valori di riferimento}

L'elettroforesi separa le proteine plasmatiche in bande caratteristiche. La
concentrazione totale normale è 6--8\,g/dl. Ciascuna zona elettroforetica
raggruppa proteine distinte:

\begin{itemize}
  \item \textbf{Zona $\alpha_1$}: $\alpha_1$-antitripsina, $\alpha_1$-glicoproteina
    acida.
  \item \textbf{Zona $\alpha_2$}: aptoglobina, $\alpha_2$-macroglobulina,
    $\alpha_2$-antiplasmina.
  \item \textbf{Zona $\beta$}: transferrina, lipoproteine a bassa densità (LDL),
    complemento C3.
  \item \textbf{Zona $\gamma$}: immunoglobuline (IgA, IgD, IgE, IgG, IgM).
\end{itemize}

Le \textbf{bande monoclonali} (picco anomalo singolo, di norma in zona
$\gamma$) sono indice di patologie quali il mieloma multiplo.

\begin{table}[htbp]
  \centering
  \small
  \renewcommand{\arraystretch}{1.3}
  \begin{tabularx}{\textwidth}{@{} l c X X @{}}
    \toprule
    \textbf{Proteina} & \textbf{g/dl} & \textbf{Aumenta in} & \textbf{Diminuisce in} \\
    \midrule
    Albumina
      & 3,6--4,9
      & Disidratazione, vomito, diarrea, sudorazione eccessiva
      & Malnutrizione, malattie renali ed epatiche, infiammazioni,
        ipertiroidismo, gravidanza \\
    $\alpha_1$-globuline
      & 0,2--0,4
      & Malattie infiammatorie croniche, infettive, infarto, gravidanza
      & Malattie epatiche gravi, enfisema congenito (ereditario), malattie renali \\
    $\alpha_2$-globuline
      & 0,4--0,8
      & Malattie renali e infiammatorie, infezioni, diabete, sindrome di Down
      & Malattie epatiche gravi, diabete, emolisi, artrite reumatoide \\
    $\beta$-globuline
      & 0,6--1,0
      & Anemia sideropenica, mieloma, ipercolesterolemia, gravidanza
      & Malnutrizione, cirrosi \\
    $\gamma$-globuline
      & 0,9--1,4
      & Gammopatie (MGUS), mieloma multiplo, epatopatie croniche, infezioni
      & Malattie ereditarie del sistema immunitario \\
    \bottomrule
  \end{tabularx}
  \caption{Frazioni proteiche plasmatiche: valori normali e principali cause di variazione.}
  \label{tab:proteine-plasmatiche}
\end{table}

\subsection{Albumina}

L'\textbf{albumina} è la proteina plasmatica più abbondante. Funziona come
proteina di trasporto a bassa specificità e amplissima capacità: lega farmaci,
bilirubina, acidi grassi, metalli e ormoni. Il suo dosaggio è un ottimo indice
della \textbf{funzionalità epatica} e la sua concentrazione diminuisce nei
processi infiammatori. L'\textbf{analbuminemia} (assenza o grave riduzione di
albumina per deficit di sintesi) non presenta una sintomatologia clinica
particolarmente marcata grazie alla compensazione delle altre frazioni.

% ---------------------------------------------------------------------------
\section{Gli eritrociti}

\subsection{Morfologia}

Gli \textbf{eritrociti} (globuli rossi, RBC) sono cellule anucleate prive di
organuli citoplasmatici. Hanno una caratteristica forma a \textbf{lente
biconcava}: diametro di 7,5\,\textmu m, spessore di circa 2\,\textmu m (1\,\textmu m
al centro). Questa conformazione massimizza il rapporto superficie/volume,
ottimizzando lo scambio di gas. Il citoplasma è occupato
dall'\textbf{emoglobina} e ricco di enzimi solubili, tra cui l'\textbf{anidrasi
carbonica}, essenziale per la formazione dello ione bicarbonato che tampona il
pH del sangue.

Il numero normale è di \textbf{5\,milioni/mm\textsuperscript{3}} nel maschio e
\textbf{4,5\,milioni/mm\textsuperscript{3}} nella femmina. La vita media è di
\textbf{120 giorni}: al termine, la membrana espone oligosaccaridi che li
rendono bersaglio dei macrofagi di milza, midollo osseo e fegato.

La \textbf{membrana plasmatica} è composta per il 50\,\% da proteine, il
40\,\% da lipidi e il 10\,\% da carboidrati; la maggior parte delle proteine è
intrinseca. Il citoscheletro — formato da spettrina, actina, anchirina e
proteine di banda 3, 4.1, 4.2 e 4.9 — costituisce un guscio unito alla
membrana in molti punti, distinguendo l'eritrocita dalle altre cellule e
conferendogli la flessibilità necessaria per percorrere i capillari,
dove viaggia impilato con gli altri eritrociti.

\subsection{Emoglobina}

L'\textbf{emoglobina} (Hb) è una proteina tetramerica (p.m. 68\,000\,Da)
formata da due catene $\alpha$ e due catene $\beta$, ciascuna legata a un
\textbf{gruppo eme} contenente Fe\textsuperscript{++}. Ogni atomo di ferro lega
reversibilmente una molecola di O\textsubscript{2}:

\begin{itemize}
  \item \textbf{HbO\textsubscript{2}} (ossiemoglobina): sangue arterioso,
    colore rosso vivo.
  \item \textbf{Hb} (deossiemoglobina): sangue venoso, colore blu-violaceo.
\end{itemize}

L'\textbf{emoglobina fetale} (HbF, catene $\alpha_2\gamma_2$) ha un'affinità
per l'O\textsubscript{2} superiore a quella dell'Hb adulta, favorendo il
trasferimento di ossigeno dalla madre al feto.

\subsection{Metabolismo del ferro}

Circa i 2/3 del ferro corporeo (2\,g nella donna, 5\,g nell'uomo) sono
legati all'Hb; 1/4 è in deposito come \textbf{ferritina}; la quota restante è
ferro funzionale (mioglobina, enzimi). Le perdite quotidiane sono di
2\,mg/die nella donna e 1\,mg/die nell'uomo. L'assorbimento avviene
prevalentemente nel \textbf{duodeno}, esclusivamente nella forma
Fe\textsuperscript{2+}.

\begin{itemize}
  \item \textbf{Carenza di ferro}: inibisce la sintesi di Hb. Cause principali:
    perdite di sangue, scarso apporto alimentare, aumentato fabbisogno
    (gravidanza), ridotto riciclo (infezioni croniche).
  \item \textbf{Sovraccarico di ferro}: danneggia fegato, pancreas e miocardio.
    Se la capacità di legame della transferrina è esaurita, il Fe libero è
    tossico per le cellule.
\end{itemize}

\subsection{Porfirie}

Le \textbf{porfirie} sono malattie metaboliche rare, per lo più ereditarie,
causate dal deficit di un enzima della via biosintetica dell'eme. Si accumulano
porfirine o loro precursori in vari tessuti, con escrezione nelle urine e nelle
feci; le manifestazioni colpiscono il \textbf{sistema nervoso} e la
\textbf{cute} (fotosensibilità).

\subsection{Anemie emolitiche ereditarie}

Nelle \textbf{anemie falciformi} e nella \textbf{sferocitosi} ereditaria, gli
eritrociti assumono una forma anomala che riduce la loro vita media: nonostante
l'incremento compensatorio dell'eritropoiesi, la capacità di trasporto di
O\textsubscript{2} rimane insufficiente.

L'\textbf{anemia falciforme} è causata da una mutazione puntiforme nel gene
della catena $\beta$ dell'Hb (GAG $\rightarrow$ GUG), che sostituisce l'acido
glutammico con la valina. I globuli rossi assumono una forma a falce che
ostacola il microcircolo. È più frequente in Africa centro-orientale, Sud
Arabia e area mediterranea (in particolare Italia, Spagna e Grecia). Le
terapie comprendono trasfusioni, splenectomia e induttori dell'Hb fetale.

% ---------------------------------------------------------------------------
\section{Gruppi sanguigni}

I \textbf{gruppi sanguigni} (sistema AB0) furono scoperti nel 1901 da Karl
Landsteiner. La membrana degli eritrociti espone catene di carboidrati
ereditarie che fungono da antigeni; gli anticorpi plasmatici corrispondenti
determinano le reazioni di incompatibilità trasfusionale.

\begin{itemize}
  \item \textbf{Gruppo A}: antigene A sulla membrana; anticorpi anti-B nel plasma.
  \item \textbf{Gruppo B}: antigene B; anticorpi anti-A.
  \item \textbf{Gruppo AB}: antigeni A e B; nessun anticorpo (ricettore
    universale).
  \item \textbf{Gruppo 0}: nessun antigene; anticorpi anti-A e anti-B
    (donatore universale).
\end{itemize}

Il \textbf{fattore Rh} fu isolato nel 1940 dalla scimmia \textit{Macacus
rhesus}. Gli individui \textbf{Rh+} esprimono l'antigene D sulla membrana
eritrocitaria; gli individui \textbf{Rh--} ne sono privi e possono sviluppare
anticorpi anti-D in seguito a esposizione a sangue Rh+.

% ---------------------------------------------------------------------------
\section{I leucociti}

\subsection{Generalità}

I \textbf{leucociti} (globuli bianchi) garantiscono la difesa dell'organismo
contro patogeni e sostanze estranee. Sono prodotti dal midollo osseo, presenti
a 5\,000--7\,000/mm\textsuperscript{3} e hanno vita media di circa 20 giorni.
Si dividono in:

\begin{itemize}
  \item \textbf{Granulociti}: citoplasma con granuli colorabili; nucleo
    plurilobato (\textit{polimorfonucleati}). Comprendono neutrofili, eosinofili
    e basofili (questi ultimi affiancati, come funzione, dai mastociti
    tissutali).
  \item \textbf{Agranulociti}: cellule mononucleate. Comprendono linfociti e
    monociti (questi ultimi, una volta emigrati nei tessuti, detti macrofagi).
\end{itemize}

In base alla funzione svolta, i leucociti si possono inoltre raggruppare in:
\textbf{fagociti} (neutrofili e monociti/macrofagi), \textbf{granulociti}
(basofili, neutrofili ed eosinofili), \textbf{cellule citotossiche} (alcuni
linfociti, tra cui le cellule natural killer) e \textbf{cellule che
presentano l'antigene} (monociti/macrofagi, linfociti B/plasmacellule e
cellule dendritiche).

\subsection{Granulociti neutrofili}

Rappresentano il \textbf{50--70\,\%} dei leucociti. Hanno diametro di
12--14\,\textmu m e nucleo plurilobato (da cui il sinonimo
\textit{polimorfonucleati}). Nelle donne può essere presente il
\textbf{corpo di Barr}, addensamento di cromatina corrispondente al cromosoma X
inattivo. Il citoplasma contiene granuli con enzimi lisosomiali e sostanze
battericide.

Sono i primi ad arrivare sul sito di lesione, mostrano spiccata attività
fagocitaria (responsabile del pus) e hanno vita breve ($\leq$12\,ore). Aderendo
all'endotelio tramite integrine e ICAM-1 (indotto da IL-1 e TNF), attraversano
la parete vascolare per diapedesi, fagocitano i batteri e rilasciano
\textbf{leucotrieni} per amplificare la risposta infiammatoria.

\subsection{Granulociti eosinofili}

Rappresentano meno del \textbf{4\,\%} dei leucociti. Contengono granuli
specifici e granuli azzurrofili: al microscopio elettronico i granuli
specifici mostrano una parte interna più densa (agenti proteici
antiparassitari e una neurotossina) e una parte esterna meno densa; i granuli
azzurrofili sono lisosomi che idrolizzano complessi antigene-anticorpo e
parassiti fagocitati. Svolgono un ruolo chiave nelle risposte
anti-parassitarie e allergiche.

\subsection{Granulociti basofili}

Meno dell'\textbf{1\,\%} dei leucociti. Il nucleo a S è mascherato da numerosi
granuli. Esprimono sulla membrana recettori per le IgE e mediano le
\textbf{reazioni allergiche} e i fenomeni di ipersensibilità immediata. Vita:
da 2\,ore a pochi giorni.

\subsection{Linfociti}

Rappresentano il \textbf{20--25\,\%} dei leucociti. Sono piccoli (8--10\,\textmu m),
con nucleo grande eccentrico e scarso citoplasma. A differenza degli altri
leucociti, sono cellule longeve e non terminali: possono trasformarsi in
\textbf{linfoblasti} e acquisire nuovi programmi funzionali dopo l'incontro con
l'antigene. Si distinguono in \textbf{linfociti T} e \textbf{linfociti B} in
base a marcatori di superficie, e costituiscono il sistema dell'immunità
specifica. Al microscopio elettronico si nota poco citoplasma, con pochi
mitocondri e numerosi ribosomi liberi, oltre ad alcuni granuli azzurrofili.

\subsection{Monociti}

Rappresentano l'\textbf{1--6\,\%} dei leucociti. Sono agranulociti con nucleo
rotondo o reniforme; una volta emigrati nei tessuti si differenziano in
\textbf{macrofagi}, cellule con elevata capacità fagocitaria e ruolo centrale
nella presentazione dell'antigene al sistema immunitario.

\subsection{Cellule dendritiche}

Dette anche \textbf{cellule di Langerhans} o \textbf{cellule velate}, sono
agranulociti specializzati nel riconoscimento dei patogeni e
nell'attivazione delle altre cellule del sistema immunitario tramite
presentazione dell'antigene.

% ---------------------------------------------------------------------------
\section{Piastrine ed emostasi}

Le \textbf{piastrine} (trombociti) sono frammenti citoplasmatici anucleati
prodotti nel midollo osseo per frammentazione dei \textbf{megacariociti},
cellule giganti caratterizzate da copie multiple di DNA nel nucleo
(poliploidia): dai margini di ciascun megacariocita si staccano decine di
migliaia di piastrine dall'aspetto di dischi irregolari. Le piastrine hanno
diametro di 2--4\,\textmu m, concentrazione di
200\,000--400\,000/mm\textsuperscript{3} e vita media di 8--10 giorni.

L'\textbf{emostasi} è il processo sequenziale che impedisce la perdita di
sangue dai vasi lesi, mantenendone l'integrità e la fluidità. La prima fase
(\textit{emostasi primaria}) è affidata all'aggregazione piastrinica nel sito
di lesione. Si definisce \textbf{trombocitosi} l'aumento e
\textbf{trombocitopenia} la riduzione del numero di piastrine rispetto ai
valori normali.
